\begin{figure}[H]
\centering
\begingroup
\small
\setlength{\tabcolsep}{8pt}
\renewcommand{\arraystretch}{1.25}
\begin{tabular}{p{0.82\textwidth}}
\toprule
\textbf{Archivo ICOCIV del DANE}: periodos reales y tablas jerárquicas. \\
\(\downarrow\) \\
\textbf{Validación de calidad de datos}: continuidad, orden, faltantes, duplicados y longitud mínima. \\
\(\downarrow\) \\
\textbf{Construcción de serie histórica}: índice mensual asociado a la selección técnica. \\
\(\downarrow\) \\
\textbf{Ajuste de modelos interpretables}: lineal, polinómico, logarítmico y robusto. \\
\(\downarrow\) \\
\textbf{Diagnóstico estadístico}: residuos, AIC/AICc, \(R^2\) ajustado, Durbin-Watson y Jarque-Bera. \\
\(\downarrow\) \\
\textbf{Backtesting temporal}: rolling forecast origin y métricas MAE, RMSE y MAPE. \\
\(\downarrow\) \\
\textbf{Proyección e informe técnico}: resultado, intervalo, trazabilidad y justificación metodológica. \\
\bottomrule
\end{tabular}
\endgroup
\par\footnotesize\textit{Fuente: Elaboración propia. Diagrama de lógica metodológica interna de la aplicación, construido con base en la validación temporal y los criterios documentados en el informe.}\normalsize
\caption{Flujo metodológico general del software ICOCIV.}
\label{fig:flujo_metodologico}
\end{figure}
